% New analytic increment only. Prior files are not modified.
\documentclass[11pt]{article}
\usepackage[T1]{fontenc}
\usepackage{lmodern,amsmath,amssymb,amsthm,mathtools,microtype,booktabs,xurl}
\usepackage[margin=.9in]{geometry}
\usepackage[hidelinks]{hyperref}
\setlength{\emergencystretch}{3em}
\allowdisplaybreaks
\numberwithin{equation}{section}
\newtheorem{proposition}{Proposition}[section]
\newtheorem{lemma}[proposition]{Lemma}
\newtheorem{corollary}[proposition]{Corollary}
\theoremstyle{remark}\newtheorem{remark}[proposition]{Remark}
\newcommand{\E}{\mathbb E}
\newcommand{\dd}{\,d}
\newcommand{\one}{\mathbf 1}
\newcommand{\Sm}{\mathcal S}
\newcommand{\Qthree}{\mathcal Q_3}
\title{Q3 strip budgets: pre-entry estimates and two stopping obstructions}
\author{Analytical increment for independent review}
\date{Version 1; October 22, 2026}
\begin{document}
\maketitle
\begin{abstract}
The S1 results imply the previously unstated pre-entry cross-residual
bound and a vanishing pre-entry signed gate budget. They yield a global
sine envelope from the $\pi$ axis with the necessary initial
$q$-regularization, and an explicit total-clock bound.
The proposed unregularized absolute radial gain is false for the
initialized original flow, already at the axis label. That stops Part A.
For Part B, neither remaining on the lower side nor uniform chart entry
of the whole $3\pi/2$ strip is true. A positive-width outer substrip stays
frozen and has a negligible mass budget; the remaining recruitment
requires an arrival-weighted estimate not supplied by the cited entry
theorems. No complete strip budget or weak-passage assembly is claimed.
\end{abstract}

\section{Setting, sources and scope}
All labels introduced here have prefix \texttt{an03q3sbv1:}.
We use the original balanced initialization and residual equations
\[
 U(0,b)=W(0,b)=q^5(\cos b,\sin b),\qquad
 U=\sqrt m\,a_\theta,\quad W=\sqrt m\,a_\psi,
\]
\[
 \dd\lambda_0=\dd b/(2\pi),\qquad
 f(X)=\int W(U\cdot X)_+\dd\lambda_0 .
\]
Here $P=(P_1+P_2)/2$, $P_1=N(e_1,q^2I)$,
$P_2=N(\rho e_2,q^2I)$, $\rho\in[13/20,7/10]$,
$e_i=X_i-f_i$, and
$R_{ij}(\theta)=\E_P[e_iX_j\one_{\{a_\theta\cdot X>0\}}]$.
The letter $b$ is an initial label, not a current angle.
Set
\begin{align}
 \lambda&=\rho^2,\quad Q=\lambda+q^2,\quad
 \omega=(1-\lambda)/2,\quad \ell=\log(1/q),\nonumber\\
 c_m&=2/\omega+1/2,\quad t_m=c_m\ell,\quad
 u=10-c_m\ge169/102,\nonumber\\
 t_0&=t_h=10\ell-\log A_q+O(q^2\ell^2+q^u),\qquad
 0<c_A\le A_q\le C_A .                         \label{an03q3sbv1:times}
\end{align}
The symbol $c_m$ is a time coefficient; $c(t)=Q^{-1}\E_{P_2}[f_2X_2]$
is the response. Define
\[
 I_{\rm tot}(t)=Q\int_0^t(1-c),\qquad
 \beta=\lambda/Q,\quad k=10(1-\lambda),\quad
 N=|\sin\theta|+|\sin\psi|,\quad
 N_-=\E_{P_2}[(f_2X_2)_-].
\]
Use the proved resident root $a_\rho=\rho K(\rho a_\rho)$ and put
$P_\rho=\Phi(\rho a_\rho)$, $d=(1-\lambda P_\rho)/2$,
$s=d/\omega$ and $\alpha=\lambda P_\rho/2$.
For the two strips use exactly GESG's constants:
\begin{equation}
 \begin{gathered}
 k_G=169/3200,\quad \Lambda_q=\sqrt{3\ell/k_G},\quad
 w_q=\arcsin(q\Lambda_q),\\
 \Sm_\pi=[\pi,\pi+2w_q),\qquad
 \Sm_{3\pi/2}=(3\pi/2-2w_q,3\pi/2].
 \end{gathered}
 \label{an03q3sbv1:strips}
\end{equation}

On $[t_0,t_*]$ retain the stopped hypotheses of S2CP/CFA, including
the instantaneous core hypotheses, cross residual $M_{\rm res}q$,
$\int r<B$, duration $t_*-t_0\le L_T\ell$, and the S2 joint secondary
guard or its proved improvement. Whenever its gate or clock conclusion
is used, its retained outside supremum is still a stopped input; it is
not proved here for the two strips. We do not assume CFA's additional
integrated strip input or invoke its conditional passage conclusion.
In particular $t_*\le C_{\rm dur}\ell$.

Source abbreviations are those of roadmap v1.13. The specific inputs
are OWPT's early mass bound; MOMC Theorem main, its mass, chart and
left-ancestry conclusions; Q2EC equations early, match and
coordinatetransport; PROFILE's exact target coordinate equation;
GESG's sine, signed and interiorspeed results; S2CP's gate lemma and
stopped clock; and TAILREP's gain-integration and tail-moment lemmas.
The review file records repository blob hashes. All arguments below
are analytic.

\section{Part A: estimates that do close before the radial step}
\begin{lemma}[Pre-entry cross residual and response budget]
\label{an03q3sbv1:pre}
For the original flow, without a new pre-entry guard,
\begin{align}
 \sup_{0\le t\le t_0}
   \big(\|e_1\|_{P_2}+\|e_2\|_{P_1}\big)&\le C_{\rm pre}q,
                                                   \label{an03q3sbv1:crosspre}\\
 \sup_{0\le t\le t_0}(|c|+N_-)
   +\int_0^{t_0}(|c|+N_-)\dd t
   &\le C(q^u+q^2\ell^2)=o(1).                     \label{an03q3sbv1:gatepre}
\end{align}
The constants depend only on the model box and the proved S1 constants,
not on LCRC's guards.
\end{lemma}
\begin{proof}
Let $S=\int m\dd\lambda_0$. OWPT gives
$S(t)\le\overline S(t)=q^{10}e^{(1+2q^2)t}$.
Thus on $[0,t_m]$ both $\sup S$ and $\int S$ are $O(q^u)=o(q)$.
The elementary bound $|f(X)|\le S|X|$ gives
$\|f_i\|_{P_j}\le CS$ and
$\E_{P_2}|f_2X_2|\le CS$.
The cross target coordinates have norm $q$.

On $[t_m,t_0]$, split into MOMC's right-half strong ancestry
$\mathcal R$ and its entire complementary ancestry. Write their masses
as $M$ and $\mu$. MOMC gives
\[
 M\le1/2,\quad \int_{t_m}^{t_0}M\le3/2,\quad
 \mu\le Cq^2\ell^2M,\quad
 |\theta|,|\psi|\le Rq\ \text{on }\mathcal R,\quad R=79/51.
\]
Consequently
\[
 \|(f_{\mathcal R})_1\|_{P_2}
 +\|(f_{\mathcal R})_2\|_{P_1}\le C_RqM,\qquad
 \|(f_{\mathcal R^c})_i\|_{P_j}\le C\mu.
\]
For the sharper response estimate, write $X=(qZ,V)$ under $P_2$,
where $V=\rho+qZ_2$. Each strong donor has
$|W_2|\le Rq\sqrt m$ and
$|U\cdot X|\le C_Rq\sqrt m(|Z|+|V|)$, whence
\[
 \E_{P_2}|(f_{\mathcal R})_2X_2|\le C_Rq^2M,\qquad
 \E_{P_2}|(f_{\mathcal R^c})_2X_2|\le C\mu .
\]
Since $Q$ is bounded below, these inequalities prove both claims.
They include all left ancestry and do not use an entry theorem for
either Q3 strip.
\end{proof}

\begin{lemma}[Local $\pi$-axis field and the signed excess]
\label{an03q3sbv1:local}
Fix $c_*>0$. Whenever $\cos\theta\le-c_*$, loss monotonicity alone
gives $|R_{11}|+|R_{21}|\le C_{c_*}q$. A bound
$\|e_1\|_{P_2}\le Cq$ additionally gives $|R_{12}|\le Cq$.
Before $t_0$ this bound follows from Lemma~\ref{an03q3sbv1:pre}.
For every receiver, before and after $t_0$, the cross entries satisfy
\[
 |R_{12}|+|R_{21}|\le a_*q,\qquad
 a_*=C(1+\max(C_{\rm pre},M_{\rm res})).
\]
Put $X(t)=\sup_\theta(|R_{22}(\theta)|-Q(1-c)/2)_+$.
For small $q$, $1-c>0$ on $[0,t_*]$ and
\begin{align}
 X(t)&\le \tfrac12 C_{\rm cross}q^2
               +\tfrac12(Qc+N_-),                 \label{an03q3sbv1:X}\\
 \int_0^{t_0}X&=o(1),\qquad
 \int_0^{t_*}X\le \frac{5e}{9}\frac{E_g}{1-c_b}+o(1).
                                                     \label{an03q3sbv1:Xint}
\end{align}
Here $C_{\rm cross}=C_{\rm pre}$ before $t_0$, and can be replaced
by a fixed multiple of $M_{\rm res}$ afterwards.
\end{lemma}
\begin{proof}
The $P_1$ gate has probability at most $e^{-c_*^2/(2q^2)}$.
Loss monotonicity and a fourth-moment bound make its selected moments
exponentially small. On $P_2$, $X_1$ has norm $q$, which pays
$R_{11},R_{21}$ using only the loss bounds. In $R_{12}$, $X_2$ has
order-one norm; this is precisely where the cross-residual bound is
needed. Clusterwise Cauchy--Schwarz also proves the global cross-entry
bound.

GESG's exact signed identities imply
\[
 \begin{split}
 R_{22}-Q(1-c)/2
   &\le C_{\rm cross}q^2/2+(Qc+N_-)/2,\\
 -R_{22}-Q(1-c)/2
   &\le C_{\rm cross}q^2/2-Q(1-2c)/2+N_-/2.
 \end{split}
\]
Here $Qc+N_-=\E_{P_2}[(f_2X_2)_+]\ge0$.
Since $c\le1$, the second right side is no larger than the first.
This proves \eqref{an03q3sbv1:X}, even if $c$ is slightly negative
before $t_0$; pre-entry nonnegativity was not assumed.
Use \eqref{an03q3sbv1:gatepre}, $q^2t_*=o(1)$ and S2CP's gate
lemma to obtain \eqref{an03q3sbv1:Xint}.
\end{proof}

For completeness, the local drift proposal has the following precise
validity. In the fixed near-$\pi$ chart
$\theta=\pi+qx,\ \psi=\pi+qy$, $|x|+|y|\le L$, the original field on
$[0,t_m]$ is
\begin{equation}
 \begin{split}
 2x'&=-\rho\phi(\rho x)+\lambda y\Phi(-\rho x)
                      +O_L(q^2+S/q),\\
 2y'&=-\rho K(-\rho x)+O_L(q^2+S/q),\qquad
 K(z)=\phi(z)+z\Phi(z).
 \end{split}                                      \label{an03q3sbv1:pifield}
\end{equation}
Indeed the $P_1$ gate is exponentially small. Under $P_2$, its leading
gate is $Z+xV<0$, giving
$R_{12}=R_{21}=-q\rho\phi(\rho x)/2+O_L(q^3)$ and
$R_{22}=\lambda\Phi(-\rho x)/2+O_L(q^2)$ for the target field.
Insert these in the exact two angular equations. The output part
costs $C_LS/q$ after scaling, as in PROFILE Proposition 4.1.
On the aligned ray $y=x$, both leading velocities equal
$-\rho K(-\rho x)/2<0$: the drift is toward Q2.
Early comparison gives $|x-y|=O(S/q)=o(1)$ on this fixed chart, so
the sign persists there for small $q$. This is not a uniform sign
estimate on $|x|\le2\Lambda_q$: there the negative Gaussian drive
can be arbitrarily small. No such growing-chart assertion is used.

\begin{proposition}[The regularized sine envelope from initialization]
\label{an03q3sbv1:envelope}
For every label in $\Sm_\pi$ and $0\le t\le t_*$,
\begin{equation}
 N(t)\le \min\{2,C_{\rm env}(N(0)+q)
                         e^{\beta I_{\rm tot}(t)/2}\}.
 \label{an03q3sbv1:Nenv}
\end{equation}
One may take, enlarging fixed constants,
\[
 a_0=\frac{\lambda_{\min}}2\min(1/2,1-c_b),\quad
 G=\frac{5eE_g}{9(1-c_b)},\quad
 C_{\rm env}=4e^{G+1}\max(1,a_*/a_0).
\]
This constant has no dependence on a strip profile or on $q$.
Moreover, at every stopped endpoint,
\begin{equation}
 e^{I_{\rm tot}(t_*)}
   \le C E_g A_q^{-q^2}q^{-10-10q^2}.
 \label{an03q3sbv1:totalclock}
\end{equation}
The constant in this last inequality is the fixed stopped-clock
constant with the S1 seed constants.
\end{proposition}
\begin{proof}
Stop first when $N=1/2$. Before this event, continuity from the
$\pi$ axis ensures $\cos\theta\le-\sqrt3/2$. GESG's matrix-value
sine inequality and the local estimate give, in upper-Dini form,
\[
 D^+N\le\{\beta I_{\rm tot}'/2+X+Cq+Cq^2\}N+a_*q.
\]
The replacement of $I_{\rm tot}'/2$ by $\beta I_{\rm tot}'/2$
costs $q^2(1-c)/2$, whose integral is $o(1)$. Also
$\beta I_{\rm tot}'/2\ge a_0$.
Duhamel and \eqref{an03q3sbv1:Xint} give
\[
 N(t)\le e^{G+1}e^{\beta I_{\rm tot}(t)/2}
                         \{N(0)+a_*q/a_0\}.
\]
If the stopping event occurs at $\tau$, its right side at $\tau$
is at least $1/2$. Since $I_{\rm tot}$ increases and $N\le2$
always, multiplying by four extends the estimate to all subsequent
times, including re-entries. This argument never uses a pre-$t_0$
integral of the global strong residual after leaving the neighborhood.

By \eqref{an03q3sbv1:gatepre} and \eqref{an03q3sbv1:times},
\[
 e^{I_{\rm tot}(t_0)}
      =A_q^{-Q}q^{-10Q}e^{o(1)}.
\]
S2CP's stopped clock, using the Q2EC/SEP seed
$H(t_0)\asymp A_q^{-\lambda}q^k$, gives
$e^{I_{\rm tot}(t_*)-I_{\rm tot}(t_0)}
 \le C E_g A_q^\lambda q^{-k}$.
Multiply; $10Q+k=10+10q^2$ and $\lambda-Q=-q^2$.
\end{proof}

\section{Part A stopping point: the proposed radial gain is false}
\begin{proposition}[An initialized-flow obstruction]
\label{an03q3sbv1:radialfailure}
There is no fixed $C$ for the proposed uniform estimate
\begin{equation}
 m(t,b)\le C m(0,b)
    \max\{1,CN(0,b)^2e^{\beta I_{\rm tot}(t)}\}^{1/\beta}
 \qquad (b\in\Sm_\pi).
 \label{an03q3sbv1:falsegain}
\end{equation}
It already fails at $b=\pi$, $t=t_m$.
The failure is an unbounded factor as $q\downarrow0$, not a
fixed-constant accuracy issue.
\end{proposition}
\begin{proof}
The exact aligned target-only flow from $b=\pi$ starts at
$\widehat U(0)=q^5(-1,0)$ and satisfies PROFILE's equation
\[
 \widehat U_2'
   =\frac{\rho q|\widehat U|}{2}
          K\left(\frac{\rho\widehat U_2}{q|\widehat U|}\right)
       +d_q\widehat U_2,\qquad 0\le d_q\le q^2.
\]
Its norm is nondecreasing, because
$(|\widehat U|^2)'=2\E_P[(\widehat U\cdot X)_+^2]\ge0$.
Also $\widehat U_2\ge0$, $K(z)\ge K(0)$ for $z\ge0$, and
$K(z)\ge z$ there. It follows that
\[
 \widehat U_2(1)\ge c q^6,\qquad
 \widehat U_2(t)\ge c q^6 e^{\lambda(t-1)/2}\quad(t\ge1).
\]
Thus
\begin{equation}
 \frac{\widehat m(t_m,\pi)}{q^{10}}
       \ge c q^{2-\lambda c_m}\longrightarrow\infty,\qquad
 \lambda c_m-2\ge\frac{210239}{184800}>0.
 \label{an03q3sbv1:seedcounter}
\end{equation}
The rational lower bound follows by monotonicity in $\lambda$ and
substitution of $\lambda=169/400$; no numerical premise is used.
Q2EC's equation early, explicitly valid for every label, gives
$|\log(m/\widehat m)|(t_m,\pi)\le Cq^u$.
But $N(0,\pi)=0$, so \eqref{an03q3sbv1:falsegain} would give
$m(t_m,\pi)\le Cq^{10}$, contradicting
\eqref{an03q3sbv1:seedcounter}.

The issue is not an irrelevant isolated endpoint. For any proposed
fixed constants and sufficiently small fixed $q$, continuous
dependence on $b$ preserves the large mass ratio on an interval
immediately to the right of $\pi$. Shrink this interval so that
$CN(0,b)^2e^{\beta I_{\rm tot}(t_m)}<1$ throughout it.
Then the same inequality fails on a positive-measure subset of Q3.
\end{proof}

The obstruction also persists if the proposed bound is requested
only at $t_0$. At the common boundary label $b=\pi$, Q2EC's band
matching and two-sided second-coordinate transport apply. This label
is in its clock cohort; the regularized RCET dichotomy and Q2EC
transport keep it off the fixed strong chart through $t_0$.
Consequently $U_2(t_0,\pi)\ge c\widehat U_2(t_0,\pi)$ and
\begin{equation}
 \frac{m(t_0,\pi)}{m(0,\pi)}
       \ge c A_q^{-\lambda}q^{2-10\lambda}\longrightarrow\infty.
 \label{an03q3sbv1:endpointcounter}
\end{equation}
Here $10\lambda-2\ge89/40>0$. This use of Q2EC concerns its
included Q2 boundary label, not an extension of its theorem to Q3.

The missing $q$ seed agrees with RCET's explicit $(s_0+q)$,
$(c_0+q)$ regularization. A repaired candidate must at least replace
$N(0)^2$ by $(N(0)+q)^2$ (or a comparable regularization).
We do not assert the repaired radial estimate.

For clarity about the remaining radial issue, the exact identity is
\[
 (\log m)'=2R_{11}\cos\theta\cos\psi
       +2R_{22}\sin\theta\sin\psi
       +2R_{12}\cos\psi\sin\theta
       +2R_{21}\sin\psi\cos\theta .
\]
DECG's projected-margin estimate, including its lower-input-side
case, and the cross bounds above give
\begin{equation}
 (\log m)'\le2|R_{11}\cos\theta\cos\psi|
       +(I_{\rm tot}'+2X)(\sin\theta\sin\psi)_+
       +C(1+\max(C_{\rm pre},M_{\rm res}))q .
 \label{an03q3sbv1:radialvalid}
\end{equation}
For negative input sine, its possible positive weak contribution is
$O(q)+O(M_{\rm res}q^2)$, as in GESG lowernogain; the $O(q)$ term
must not be omitted. Before $t_0$, use $C_{\rm pre}$ in that estimate.
This differential bound and TAILREP do not remove the inhomogeneous
seed in \eqref{an03q3sbv1:Nenv}.

After departure from a fixed $\pi$ neighborhood, the small local
$R_{11}$ estimate also no longer follows from a closed $P_1$ gate.
The post-$t_0$ contribution is paid by $\int r<B$.
A direct absolute-value radial proof from initialization would
still need, for its remaining pre-entry times, the uniform bound
\begin{equation}
 \sup_{b\in\Sm_\pi}
 \int_{\substack{0\le t\le t_0\\ \cos\theta(t,b)>-c_*}}
 |R_{11}(\theta(t,b))\cos\theta(t,b)\cos\psi(t,b)|\dd t
 \le C_{\pi,\rm rad}.
 \label{an03q3sbv1:unpaidradial}
\end{equation}
A signed substitute sufficient for the radial upper bound would also
do. This display is an unpaid consumer inequality, not an added
hypothesis of a theorem. The envelope's trivial bound $N\le2$ after
exit does not pay it. Part A stops at the false central estimate
\eqref{an03q3sbv1:falsegain}; neither A4's radial extension nor A5's
full-strip mass target is claimed.

\section{Part B: joined labels and a paid frozen substrip}
\begin{proposition}[The lower-side shortcut fails]
\label{an03q3sbv1:lowerfailure}
It is false that every $\Sm_{3\pi/2}$ label has
$\sin\theta<0$ throughout $[t_0,t_*]$.
For each sufficiently small fixed $q$, a positive-measure interval
of this strip adjacent to $3\pi/2$ has $\sin\theta(t_0)>0$.
\end{proposition}
\begin{proof}
The label $b=3\pi/2$ is the same initialized characteristic as
$b=-\pi/2$ by periodicity. MOMC includes the latter endpoint in its
right-half ancestry. At $t_m$ its scaled displacement from the actual
reference is $O(q^{d/2})$ and that reference is $(a_\rho,a_\rho)+o(1)$,
with $a_\rho>0$. MOMC's mean and shape conclusions retain this
positive chart at $t_0$. Thus
$\theta(t_0,3\pi/2)/q=a_\rho+o(1)>0$, using the lift near zero.
Continuous dependence on the initial label, for fixed $q>0$, gives
the asserted neighboring interval. No width uniform in $q$ is
inferred from this continuity argument.
\end{proof}

\begin{proposition}[A positive-width frozen part of the same strip]
\label{an03q3sbv1:frozen}
Let the fixed initial substrip be
\[
 F_q=[3\pi/2-7w_q/4,\;3\pi/2-3w_q/2].
\]
Uniformly for $b\in F_q$ and $0\le t\le t_*$,
\begin{align}
 |\theta(t,b)-b|+|\psi(t,b)-b|&\le Cq^3\ell,\nonumber\\
 m(t,b)&=q^{10}e^{O(q^3\ell)}.                    \label{an03q3sbv1:freeze}
\end{align}
In particular
\begin{align}
 \sup_{[t_0,t_*]}\mu_{F_q}/q&\le Cq^{10}\ell^{1/2},\nonumber\\
 \int_{t_0}^{t_*}\mu_{F_q}/q&\le Cq^{10}\ell^{3/2}=o(1).
                                                        \label{an03q3sbv1:Fpaid}
\end{align}
Nevertheless these labels are outside every fixed scaled strong chart
at $t_m$ and $t_0$. Their joint scaled displacement from the actual
strong reference at $t_0$ is at least $c/q$.
\end{proposition}
\begin{proof}
Write $\delta=3\pi/2-\theta$. Until the first exit from
$w_q<\delta<3w_q$, one has, for small $q$,
$\cos\theta=-\sin\delta\le-q\Lambda_q$ and
$\sin\theta=-\cos\delta\le-q\Lambda_q$.
GESG interiorspeed is valid before $t_0$ as well as afterwards,
using loss alone. All angular and logarithmic radial speeds
are at most
$C e^{-k_G\Lambda_q^2}=Cq^3$.
The distance of each initial angle in $F_q$ from either buffer
boundary is at least $w_q/2$. The available travel
$CC_{\rm dur}q^3\ell=o(w_q)$ rules out the first exit.
Integration proves \eqref{an03q3sbv1:freeze}.

The initial angular width is $w_q/4\asymp q\ell^{1/2}$ and
$\dd\lambda_0=\dd b/(2\pi)$, which proves \eqref{an03q3sbv1:Fpaid}.
The actual strong reference at $t_0$ has both angles $O(q)$,
whereas these input and output angles remain within $O(w_q)$
of $-\pi/2$ modulo $2\pi$. Even the shortest angular distance
to that reference is bounded below by a fixed positive number.
Division by $q$ proves the last claim.
\end{proof}

Thus the proposed extension of MOMC's \emph{whole-family bounded chart}
to $\mathcal R\cup\Sm_{3\pi/2}$ is false:
\begin{equation}
 \text{there is no fixed }Z\text{ with }
 \sup_{b\in\Sm_{3\pi/2}}\xi(t_0,b)\le Z .
 \label{an03q3sbv1:falsechart}
\end{equation}
The same obstruction occurs at $t_m$. It is not a bounded-factor
loss in a comparison constant. PROFILE's right-half formulas do not
assert uniform arrival for these negative-cosine labels.

This does \emph{not} disprove a raw mass-weighted tail or the desired
mass budget: the frozen labels just exhibited have negligible mass.
The two failed premises are the lower-side shortcut and whole-strip
chart extension. The remaining joined/transition population needs the
arrival-matched alternative named in the question, which is not a
conclusion of MOMC's left-mass estimate.

\section{Exact remaining consumer and constant audit}
For clarity we record what has not been supplied, without assuming it
in a claimed full theorem. For already charted labels $J_0$ in the
unpaid part of $\Sm_{3\pi/2}$, S2CP's consumer is the raw law
\begin{equation}
 \eta_{J_0}\{\xi>v\}\le C_T\min\{1,(\epsilon_h/v)^p\},
 \quad \dd\eta_{J_0}=m(t_0,b)\dd\lambda_0,\quad
 p=3/s,\quad
 \epsilon_h=\sqrt2 A_q^d(1-M_0)^\alpha q^{10d-s},
 \label{an03q3sbv1:missingtail}
\end{equation}
together with the appropriate core cutoff and a budget for the
remaining labels. Equivalently, TAILREP consumes a bound of the form
\begin{equation}
 \int_{J_0\cap\{\xi>d_c\}}
       \max\{1,A(t)\xi^2\}^{1+\nu}\dd\eta_{J_0}
 \le C_T'\epsilon_h^p d_c^{-p}
       \max\{1,A(t)d_c^2\}^{1+\nu},\quad
 A(t)=A_0q^2e^{\beta[I_{\rm tot}(t)-I_{\rm tot}(t_0)]}.
 \label{an03q3sbv1:missingmoment}
\end{equation}
Here $d_c$ is TAILREP's fixed core cutoff. This is a mass-weighted
inequality, not a count of labels; it is not proved by
$\mu_{\rm left}(t_0)\le Cq^2\ell^2$.

An arrival-matched formulation for the unpaid population
$J=\Sm_{3\pi/2}\setminus F_q$ can be made exact as follows.
Let $\tau_b\ge t_0$ be its first joint entry into the chosen bounded
actual-reference chart, taking $\tau_b=t_0$ for labels already there
and $\tau_b=\infty$ for those not arriving. Denote its scaled
entry distance by $\xi_b$. Define
\begin{align}
 \mathcal W_J(t)={}&
   \int_{J\cap\{\tau_b>t\}}m(t,b)\dd\lambda_0(b)\nonumber\\
 &+C_m\int_{J\cap\{\tau_b\le t\}}m(\tau_b,b)
    \max\{1,A_0q^2\xi_b^2
       e^{\beta[I_{\rm tot}(t)-I_{\rm tot}(\tau_b)]}\}^{1+\nu}
          \dd\lambda_0(b).                         \label{an03q3sbv1:arrival}
\end{align}
For the second term to dominate current mass one must establish the
DECG-type gain from the actual arrival, with the remaining stopped
budgets; it cannot be assigned a fictitious $t_0$ entry.
The still-needed aggregate estimates are
\begin{equation}
 \sup_{[t_0,t_*]}\mathcal W_J(t)/q=o(1),\qquad
 \int_{t_0}^{t_*}\mathcal W_J(t)/q\dd t=o(1),
 \label{an03q3sbv1:missingarrival}
\end{equation}
or the second estimate with CFA's specified small fixed allowance.
They state the required correlation of arrival delay, entry mass and
entry displacement. Neither the right-half normalized tail nor the
unweighted left-ancestry mass gives that correlation.
No arrival theorem or additional hypothesis is asserted here.

The order-one constants in the closed results are as follows.
The loss/Gaussian constants, $k_G$, $R=79/51$, $c_A,C_A$,
$C_{\rm pre}$, and the fixed-chart coefficients in
\eqref{an03q3sbv1:pifield} are independent of LCRC guards.
$C_{\rm dur}$ uses the retained duration and the bounded $A_q$ shift.
The sine constant uses $a_*=C(1+\max(C_{\rm pre},M_{\rm res}))$,
$c_b$, and the already proved gate budget $5eE_g/[9(1-c_b)]$.
Dependence on $B$ is not needed in that sine proof; the unpaid radial
route uses $B$ only after $t_0$. The total-clock constant is inherited
from the stopped clock and S1 seed, with the explicit remaining
$A_q^{-q^2}$ in \eqref{an03q3sbv1:totalclock}.
The constants $C_T',A_0,C_m$ in the unpaid consumer are not newly
constructed or improved here. When available, their order is the
S2CP/CFA order, and the gain exponent retains
$1/\beta=1+q^2/\lambda$ (or its fixed upper bound $1+\nu$).

For $F_q$, \eqref{an03q3sbv1:Fpaid} is smaller than any fixed CFA strip
allowance for small $q$. For the \emph{whole} strips we have not supplied
CFA's bounds
\[
 \sup\mu_{\Sm}/q=o(1),\qquad
 C_{\Sm}\int\mu_{\Sm}/q
   \le\min(P_{\rm allow}/4,\overline F_F/4).
\]
Accordingly their placement in CFA's constant order is unchanged.
Part A is partial and stopped at \eqref{an03q3sbv1:falsegain}.
Part B is partial and stopped at \eqref{an03q3sbv1:falsechart}
and the unpaid arrival consumer \eqref{an03q3sbv1:missingarrival}.
There is no full strip theorem, no application of CLOCK passage,
and no joint weak-passage bootstrap in this increment.
\end{document}

